\scanpage{159}
Catégorie $\unsure{\mathcal C|\mathscr C}$, sous-catégorie $\unsure{\mathscr P|\mathcal P}$\deleted{\unsure{,}} (pleine), morphismes ``non ramifiés''

\uline{1°} Dans $\mathcal C$, les sommes finies existent, ainsi que les produits fibrés.

\uline{4°} Dans $\mathcal C$, le \deleted{produit} \inserted{composé} de deux morphismes non ramifiés, \inserted{les morphismes identiques} ou l'image
inverse d'un morphisme non ramifié, sont non ramifiés, \deleted{\unsure{id. pour}}

\uline{5°} \deleted{\unsure{si}} $f:X\to Y$ est non ramifié $\Longleftrightarrow$ les $\unsure{f_i|f}\vert X_i$ ($X_i$ les comp. \unsure{connexes|conn.}) \deleted{\unsure{la somme|les sommes} $\coprod \unsure{X_i'}\to Y$} non ramifiés

\uline{2°} Dans $\mathcal C$, il y a des composantes connexes.

\uline{3°} Les \deleted{\unsure{pts|points}} $\xi\in\mathscr P$ sont connexes (donc $\operatorname{Hom}(\xi,X\amalg Y)=\operatorname{Hom}(\xi,X)\amalg\operatorname{Hom}(\xi,Y)$]

\uline{6°} Pour tout $f:X\to Y$ non ramifié, \inserted{$Y$ connexe,} et $\xi\in\mathscr P$, $\operatorname{Hom}(\xi,X)\mathrel{\unsure{-|\to}}\operatorname{Hom}(\xi,Y)$ est
\uline{surjectif}

\uline{7°} Si \deleted{\unsure{$X$ et}} $Y$ sont connexes, et la fibre de $\operatorname{Hom}(\xi,X)\to\operatorname{Hom}(\xi,Y)$ est
réduite à un élément, alors \deleted{\unsure{$X$}} $f$ est un \uline{isom}
(on suppose $f$ non ramifié, $\xi\in\mathscr P$)

\uline{8°} Si $f:X\to Y$ non ramifié, $Y$ connexe, alors \deleted{pour} toute
\deleted{\unsure{les}} section $g:Y\to X$ est un \uline{isom.} de $Y$ sur une
composante connexe de $X$.

\uline{9°} Soit \deleted{$G$ \inserted{\deleted{groupe fini}} opère dans $X$ ``sans poi} $f:X\to Y$ \uline{ramifié}, et
$X$ connexe, \deleted{\unsure{$\xi\in$}} $e\in\operatorname{Hom}(\xi,Y)=\unsure{:|\,}Y^\xi$, supposons $G=\operatorname{Aut}_Y(X)$
opère transitiv\unsure{ement|emt} dans $(f^\xi)^{-1}(e)$. Alors pour tout
sous-groupe $H$ de $G$, $X/H$ existe \inserted{dans $\mathcal C$} et est un
revêtement non ramifié de $Y$. \deleted{et $X$ est un \unsure{rev. princ. pour} $H$,}
i.e. deux élts de $\operatorname{Hom}(\xi,X)$ \unsure{au dessus de} $e$
ayant \unsure{la même image} par \deleted{\unsure{$X$}} $\mathrel{-} X/H$ sont \unsure{congrus|conjugués} mod $H$.

Sous ces conditions, pour $x\in\operatorname{Hom}(\xi,X)$ on définit
$\pi_1(X,x)$. Il ne change pas \deleted{\unsure{si}} en remplaçant $\xi$ par $\xi'$ tel que
$\xi'\mathrel{-}\xi$ soit donné. C'est un foncteur \uline{covariant} en $(X,x)$

\diagram{\begin{tikzcd}[ampersand replacement=\&]
\xi \arrow[r,dash,"x"] \arrow[d] \& X \arrow[d]\\
\xi' \arrow[r,dash,"x'"] \& X'
\end{tikzcd}}
définissant $\pi_1(X,x)\mathrel{-}\pi_1(X',x')$.

\uline{Etc.}
% archive sheet number: 158
